% !TeX root = heights-and-equi-distribution-of-small-points.tex
\section{Heights and equi-distribution of small points}

The main references are \cite{Yuan-Guo-2025-FirstCourseArakelov}.

We also refer to \cite[Appendices~A.5--A.6 and \S 5.4]{Yuan-Zhang-2026-AdelicLineBundles} and \cite[Chapter~3]{Chambert-Loir-2012-ArakelovGeometryHeights}.

We will use the language of Berkovich space, for reference, see \cite{Berkovich-2012-SpectralTheoryAnalytic,Ducros-Favre-Nicaise-2015-BerkovichSpacesApplications}.

\begin{notation}[label=\labelNFM{not:berkovich-space-and-a-geometry}, status=draft, tag=0000000]
	We fix some notations that will be frequently used in this section.
	\begin{itemize}
		\item \latexdeleted{Recall that a \emph{valuation field} is a ...}
		      \begin{latexmathchange}
		      A \emph{valuation field} $(\kk,|\cdot|)$ is a field with an absolute value:
		      $|a| = 0$ if and only if $a = 0$, $|ab| = |a||b|$, and $|a+b| \leq |a|+|b|$.
		      Completeness always refers to the metric $|a-b|$.
		      \end{latexmathchange}
		      It is non-archimedean if $|a + b| \leq \max\{|a|,|b|\}$.
		      Two nontrivial absolute values define the same \emph{place} if they differ by a positive real power.
		      \Yang{}
		      \latexdeleted{To Ai: finish the notation, both $\calO_\kk$ and $\kk^\circ$}
		      \begin{latexmathchange}
		      For a non-archimedean valuation field, put
		      \[ \calO_\kk = \kk^\circ \coloneqq \{a \in \kk : |a| \leq 1\}, \qquad \frakm_\kk = \kk^{\circ\circ} \coloneqq \{a \in \kk : |a| < 1\}. \]
		      Then $\kk^\circ$ is a valuation ring with fraction field $\kk$, unique maximal ideal $\frakm_\kk$, and residue field
		      \[ \calk = \widetilde{\kk} \coloneqq \kk^\circ/\kk^{\circ\circ}, \qquad (\kk^\circ)^* = \{a \in \kk : |a| = 1\}. \]
		      These notations also apply to complete extensions such as $\scrH(x)$.
		      The ring $\calO_K$ of integers of a number field is a global object; at a finite place $v$, the corresponding local valuation ring is $\calO_{K_v} = K_v^\circ$.
		      \end{latexmathchange}
		\item Let $K$ be a number field, $M_K$ its set of places, and $M_K^f$, $M_K^\infty$ its subsets of finite and infinite places.
		      Write $K_v$ for the completion at $v$.
		      For $v \mid p$, normalize $|\cdot|_v$ to extend the usual absolute value on $\bbQ_p$, with $|p|_p = p^{-1}$ at finite places and $\bbQ_\infty = \bbR$.
		      Set $n_v \coloneqq [K_v:\bbQ_p]/[K:\bbQ]$ for $v \mid p$ and $p \in M_\bbQ$ \latexdeleted{and $n_v$} \Yang{}
		      \begin{latexmathchange}
		      Thus $n_v = 1/[K:\bbQ]$ at a real place and $n_v = 2/[K:\bbQ]$ at a complex place, using the usual complex modulus.
		      Write $\bbC_v$ for the completion of an algebraic closure of $K_v$; at an infinite place this means $\bbC$.
		      \end{latexmathchange}
		      Then we have the product formula
		      \[ \sum_{v \in M_K}n_v\log|a|_v = 0 \quad (a \in K^*). \]
		\item Unless otherwise stated, $X$ is a geometrically integral projective variety over $K$, $d = \dim X \geq 1$, and $L$ is a line bundle on $X$.
		      Write $X_{\bar{v}}^\an \coloneqq (X_{\bbC_v})^\an$ for its Berkovich analytification; see \Yang{}.
		      \begin{latexmathchange}
		      In the local measure notation below, $X_v^\an$ is also an abbreviation for $X_{\bar v}^\an$.
		      At an infinite place this is $X(\bbC)$ with its usual analytic topology.
		      Thus Galois conjugates are averaged on the space over $\bbC_v$, before taking any quotient by the local Galois group.
		      \end{latexmathchange}

		\item Let $\kk$ be a field and $R$ be a ring such that $\Frac R = \kk$.
		      Let $X$ be a variety over $\kk$.
		      \latexdeleted{A \emph{flat model} of $X$ over $R$ is}
		      \begin{latexmathchange}
		      A \emph{flat model} of $X$ over $R$ is a flat separated $R$-scheme $\calX$ of finite type, together with a specified isomorphism $\calX_\kk \cong X$.
		      Here $R$ is understood to be a domain.
		      \end{latexmathchange}
		      If $X/\kk$ is proper (resp. projective), we require that $\calX/R$ is also proper (resp. projective).
		      \Yang{}
	\end{itemize}
	\Yang{}
\end{notation}

\begin{proposition}[label=\labelNFM{prop:spreading-out-for-berkovich}, status=draft, tag=0000000]
	\latexdeleted{Let $R$ be a integral ring with $\Frac R = \kk$ and $X/\kk$ a variety.}
	\latexdeleted{Then $X$ has a flat model.}
	\latexdeleted{If $X/\kk$ is proper or projective, then $\calX/R$ can also be taken to be...}
	\latexdeleted{When $R$ is a valuation ring and $\calX/R$ is proper, we have the reduction map $X(\overline{\kk}) \to \calX_\calk$.}
	\latexdeleted{To Ai: finish the prop and give a references}
	\begin{latexmathchange}
	Let $R$ be a domain with fraction field $\kk$, and let $X/\kk$ be a variety.
	\begin{enumerate}
		\item There exists $0 \neq s \in R$ such that $X$ has a flat model of finite presentation over $R[1/s]$.
		      If $X$ is proper, respectively projective, the model can be chosen proper, respectively projective.
		\item If $R = \kk^\circ$ for a non-archimedean valuation field with a real-valued absolute value, $X$ has a flat model over $R$ itself.
		      Properness and projectivity can again be preserved.
		\item Suppose moreover that $\kk$ is complete and $\calX/R$ is a proper flat model.
		      There is a reduction map
		      \[ \frakr_\calX : X^\an \longrightarrow |\calX_\calk|. \]
		      For a classical point, this is obtained by extending the point over the valuation ring of a finite extension of $\kk$ and taking its special point.
	\end{enumerate}
	For spreading out and generic flatness, see \cite[Tags \href{https://stacks.math.columbia.edu/tag/01ZM}{01ZM}, \href{https://stacks.math.columbia.edu/tag/051S}{051S}]{Stacksprojectauthors-2026-StacksProject}.
	For the valuation-ring argument and proper extension of points, see Tags \href{https://stacks.math.columbia.edu/tag/0539}{0539} and \href{https://stacks.math.columbia.edu/tag/0BX4}{0BX4} of the same reference.
	\end{latexmathchange}
\end{proposition}
\begin{proof}
	\Yang{}
	\begin{latexmathchange}
	Finite presentation allows the equations and gluing data of $X$ to descend to $R[1/s]$.
	After enlarging $s$, properness or a chosen projective embedding also descends.
	Generic freeness on a finite affine covering then makes this model flat after a further localization.
	This explains why the assertion over a general domain is a statement over an open part of its spectrum.

	Now let $R = \kk^\circ$.
	If the valuation is trivial, take $\calX = X$.
	Otherwise choose $a$ with $0 < |a| < 1$.
	For every $b \in \kk$, one has $a^n b \in R$ for sufficiently large $n$, so $\kk = R[1/a]$.
	Consequently $X$, viewed as an $R$-scheme, is already a flat model of finite presentation, though its special fiber is empty.
	If $X$ is proper over $\kk$, apply Nagata compactification over $R$ and take the schematic closure of $X$ in the compactification.
	Its generic fiber is $X$, because a proper open subscheme of the separated generic fiber is also closed there.
	The closure is integral, hence torsion free over $R$, hence flat by Tag \href{https://stacks.math.columbia.edu/tag/0539}{0539}.
	For a projective $X$, use instead its schematic closure in $\bbP^N_R$.
	Nagata compactification is \cite[Tag \href{https://stacks.math.columbia.edu/tag/0F41}{0F41}]{Stacksprojectauthors-2026-StacksProject}.

	Finally, for $x \in X^\an$, properness extends $\Spec\scrH(x) \to X$ uniquely to $\Spec\scrH(x)^\circ \to \calX$.
	The image of the closed point lies in $\calX_\calk$ and defines $\frakr_\calX(x)$.
	For $x \in X(\overline{\kk})$, completeness gives the unique extension of the absolute value to its finite field of definition.
	The same valuative criterion gives its reduction.
	The target is the underlying special-fiber space; its points need not be $\calk$-rational.
	\end{latexmathchange}
\end{proof}















\subsection[status=draft,tag=0000000,label=\labelNFM{subsec:adic-metric-and-model-metric}]{Berkovich space, adic metric and model metric}

On this subsection, fix a non-archimedean completed valuation field $\kk$.

Let $X$ be a variety over a complete valuation field $\kk$.
Recall its Berkovich analytification.

\begin{construction}[label=\labelNFM{constr:berkovich-analytification}, status=draft, tag=0000000]
	Let $X$ be a variety over a complete valuation field $(\kk,|\cdot|)$.
	For an affine open $U = \Spec A \subset X$, its \emph{Berkovich analytification} is defined as
	\[ U^\an \coloneqq \left\{  \text{  multiplicative seminorm } \rho: A \to \bbR_{\geq 0} \text{ whose restriction on } \kk \text{ is } |\cdot| \right\}. \]
	Its topology is the weakest one for which $x \mapsto \|a\|_x$ is continuous for every $a \in A$.
	The affine spaces glue to a set $X^\an$.
	We have
	\[ X^\an = \left\{ (x,\gamma) \;\middle|\; x \in |X|, \gamma \text{ an absolute value on } \rkk(x) \text{ extending } |\cdot| \text{ on } \kk \right\}. \]
	Given $\rho \in U^\an$, we can get a scheme-theoretic point $\wp_\rho \coloneqq \Ker \rho$ and a absolute value on $\rkk(\wp_\rho) = \latexdeletedmath{A/\wp_\rho}\latexaddedmath{\Frac(A/\wp_\rho)}$ by $\bar{a} \mapsto \rho(a)$ \latexmath{and extension to fractions}.
	Conversely, given $(x,\gamma)$ and any affine neighborhood $U$ of $x$, we can define a multiplicative seminorm by $\calO_X(U) \to \rkk(x) \xrightarrow{\gamma} \bbR_{\geq 0}$.

	On $X^\an$, there is a Grothendieck topology.
	\Yang{}


	For $x \in X_v^\an$, let $\scrH(x)$ denote its completed residue field.
	\Yang{}
\end{construction}


Finally, let us recall some useful construction and properties of Berkovich space.

\begin{construction}[label=\labelNFM{constr:related-construction-in-berkovich-space}, status=draft, tag=0000000]
	Let us recall some construction; see \cite[Section 9.1.1]{Yuan-Guo-2025-FirstCourseArakelov}.
	\begin{itemize}
		\item \emph{Completed residue field.} Let $X/\kk$ be a variety and $(x,\rho) \in X^\an$.
		      The completed residue field is defined as the completion of $\rkk(x)$ with respect to $\rho$, denoted by $\scrH(x)$.
		\item \emph{Classical points.} Recall that for a finite extension $\bfl/\kk$ with $\kk$ completed valuation field, there exists a unique absolute valuation $\rho_{\bfl/\kk}$ on $\bfl$ extending the one on $\kk$.
		      Then we have a map
		      \[ \fraki: X(\overline{\kk}) \to X^\an, \qquad x \mapsto (x,\rho_{\rkk(x)/\kk}). \]
		      \latexdeleted{This map is the unique section of the forget map $X^\an \to |X|$ and hence injective.}
		      \begin{latexmathchange}
		      Geometric points with the same $\kk$-Galois orbit have the same image in $X^\an$.
		      More precisely, closed scheme points, rather than all geometric points separately, give an injection $|X|_{\mathrm{cl}} \hookrightarrow X^\an$.
		      For example, over $\bbQ_p$ all conjugates of an algebraic number give the same seminorm on $\bbQ_p[T]$.
		      Over an algebraically closed field the classical-point map itself is injective.
		      \end{latexmathchange}
		\item \emph{Functoriality}
		\item \emph{Reduction map.} Let $X/\kk$ be a projective variety with a projective integral model $\calX/\calO_\kk$.
		      For every $x = (x,\rho) \in X^\an$, by the valuation criterion, we have the diagram
		      \[
			      \begin{tikzcd}
				      \Spec \scrH(x) & \Spec \rkk(x) & \calX \\
				      \Spec \calO_{\scrH(x)} = \{\scrH(x), \calk_{\scrH(x)}\} & & \Spec \calO_\kk = \{\kk,\calk\}.
				      \arrow[from=1-1, to=1-2]
				      \arrow[from=1-2, to=1-3, "x"]
				      \arrow[from=1-1, to=2-1]
				      \arrow[from=1-3, to=2-3]
				      \arrow[from=2-1, to=2-3]
				      \arrow[from=2-1, to=1-3, gray, "\exists! \tilde{x}"]
			      \end{tikzcd}
		      \]
		      Set the reduction map as
		      \[ \frakr_\calX: X^\an \to \calX_\calk, \qquad x \mapsto \tilde{x}(\calk_{\scrH(x)}). \]
		      This reduction map $\frakr_\calX$ extends the reduction map of classical points in \cref{prop:spreading-out-for-berkovich} along $\fraki$.
	\end{itemize}



	\Yang{}
\end{construction}


Let us turn to the metric on line bundles.

\begin{definition}[label=\labelNFM{def:adic-metric}, status=draft, tag=0000000]
	Let $L$ be a line bundle on $X$.
	A \emph{continuous $v$-adic metric} on $L$ is a family of norms $\|\cdot\|_{x}$ on the one-dimensional $\scrH(x)$-vector space $L(x)$ for every $x \in X_{\bar{v}}^\an$ for any continuous local section $s \in \Gamma(L,U;\scrC)$, the function $x \mapsto \|s(x)\|_x$ is continuous;
	and is Galois-invariant, i.e., $\|\sigma(s)\|_{\latexdeletedmath{x}\latexaddedmath{\sigma(x)}} = \|s\|_x$ for all $\sigma \in \Gal(\overline{K_v}/K_v)$.
	The pair $\bar L_v = (L,\|\cdot\|_v)$ is called a \emph{metrized line bundle}.
\end{definition}

\begin{example}[label=\labelNFM{eg:standard-metrics-on-projective-space}, status=draft, tag=0000000]
	Let $\bbP^d$ with coordinates $[x_0:\cdots:x_d]$ and $L = \calO(1)$.
	The following rule gives a $v$-adic metric on $L$:
	\[ X_k \mapsto \frac{|X_k|_v}{\max_{0 \leq i \leq d}|X_i|_v}. \]
	This is called the \emph{standard metric} on $\calO(1)$.
\end{example}

Now we introduce a class of metric, which miraculously serves as a non-archimedean analoge of $\cinf$-function.

First let see the complex case, $(X,L)/\bbC$.
To determine a $\cinf$-hermitian metric $\|\cdot\|$ on $L$, we only need to choose an everywhere nonvanish $\cinf$-section $s$ with $\|s(x)\| = 1$ for all $x \in X$.
So, choose an everywhere nonvanish $\cinf$-section $s$, we can get a hermitian metric on $L$ by $\|as\| \coloneqq |a|$ for any $\cinf$-function $a$.
More generally, a $\cinf$-orthogonal frame, i.e., a collection of trivialization $\{(U_i,L|_{U_i} \cong \cinf_{X,\bbC} \cdot e_i)\}$ with transition function $e_i = f_{ij} e_j$ such that $|f_{ij}|=1$, determines a $\cinf$-hermitian metric on $L$ by $\|ae_i\| \coloneqq |a|$ for any $\cinf$-function $a$.
Hence, a $\cinf$-hermitian metric is determined by a collection of transition functions $\{f_{ij}\}$, where
\[ \bigl( f_{ij}: X \to S^1 = \{|z|=1\} \subset \bbC^* \bigr) \in \cinf(X,S^1). \]

On the non-archimedean place, we have no good $\cinf$-functions.
However, the elements of norm $1$ has some very good structure, namely, we can consider $(\kk^\circ)^*$, the multiplicative group of the ring of integers.
Hence, the question becomes, what is a suitable function class (?).
\[ \bigl( f_{ij}: X \to (\kk^\circ)^* = \{|a|=1\} \subset \kk^*\bigr) \in (?). \]
For a scheme $\calY/(\kk^\circ)$, there is a natural class of functions $\calY \to (\kk^\circ)^*$ given by $\calO_\calY^*$.
Therefore, for the variety $X/\kk$, the natural way should be that taking an integral model $\calX/(\kk^\circ)$ and by $X \to \calX \to (\kk^\circ)^*$.

\begin{construction}[label=\labelNFM{constr:model-metric}, status=draft, tag=0000000]
	Let $X$ be a projective variety over $\kk$ and $L$ line bundle on $X$.
	Given an integral model $(\calX,\calL)$, where $\calX$ is an integral model of $X$, $\calL$ a line bundle on $X$ with an isomorphism $\calL|_X \cong L$,
	then we can define a \emph{model metric} $\|\cdot\|_{(\calX,\calL)}$ by the frame of $\calL$ on $X$.
	More explicitly, let $\{(\calU_i,\calL|_{\calU_i} = \calO_{\calU_i} \cdot e_i) \}$ be a trivialization of $\calL$.
	This also gives a trivialization
	\[ \{((\calU_i \cap X \eqqcolon U_i)^{\an}, L|_{U_i}^{\an} = \calL|_{U_i}^{\an} \cong \scrO_{U_i^\an} \cdot e_i\} \]
	of $L^{\an}$ on $X^\an$.
	For any local continuous section $ae_i$, set $\|ae_i(x)\|_{(\calX,\calL)} \coloneqq |a|_x$.

	For latter use, we also need a version with rational coefficients.
	Given an integral model $(\calX,\calM)$ of $(X,L^{\otimes m})$, the \emph{model metric (with rational coefficients)} of $(\calX,\calM, m)$ is defined as
	\[ \|s\|_{(\calX,\calM,m)} \coloneqq \|s^{\otimes m}\|_{(\calX,\calM)}^{1/m}.\]
\end{construction}

\begin{example}[label=\labelNFM{eg:standard-metric-are-model-metric}, status=draft, tag=0000000]
	Let $\|\cdot\|$ be the standard metric on $(\bbP^d,\calO(1))$.

	\Yang{}
\end{example}

\begin{definition}[label=\labelNFM{def:semipositive-local-metric}, status=draft, tag=0000000]
	A continuous metric is called \emph{semipositive} if it is a uniform limit of the following metrics: model metrics for which $\calL$ is relatively nef, that is, $\deg(\calL|_C) \geq 0$ for every integral curve $C$ in the special fiber.
	\Yang{}
\end{definition}


Recall that in the archimedean case, any $\cinf_{-,\bbR}$-function can be got by $-\log \|1\|$ for some $\cinf$-metric $\|\cdot\|$ on $\scrO_X$.
By this way, we have the following define.

\begin{definition}[label=\labelNFM{def:model-function-over-na-place}, status=draft, tag=0000000]
	Let $X$ be a projective variety over $\kk$.
	A \emph{model function} on $X$ is a function $X^\an \to \bbR$ of form $- \log \|1\|_{(\calX,\calL,m)}$, where $(\calX,\calL,m)$ is a model of $(X,\calO_X)$ and $1 \in \calO_X(X)$ is the canonical section.
\end{definition}

\begin{example}[label=\labelNFM{eg:model-function}, status=draft, tag=0000000]
	\latexdeleted{To Ai: give some example and nonexample of model function}
	\begin{latexmathchange}
		Let $\pi$ be a uniformizer of $K_v$.
		Since the special fiber of a model is the divisor of $\pi$, every rational multiple of $-\log|\pi|_v$ is a constant model function.

		For a nonconstant example, let $X = \bbP^1_K$ with affine coordinate $z$ and put
		\[ \varphi(z) \coloneqq \log\frac{\max\{|\pi|_v,|z|\}}{\max\{1,|z|\}}. \]
		The expression extends continuously to $\bbP^{1,\an}_{\bbC_v}$, with value $0$ at $\infty$.
		It is a model function: on $\bbP^1_{K_v^\circ}$, take the vertical ideal $(\pi,z)$ on the chart $z \neq \infty$ and the unit ideal on the other chart.
		Its blowup makes the ideal invertible; the negative of the exceptional divisor gives $\varphi$.

		The continuous function $\varphi^2$ is \emph{not} a model function.
		Indeed, on the segment of Gauss points $x_{0,r}$ with $|\pi|_v \leq r \leq 1$,
		\[ \varphi(x_{0,r}) = \log r, \qquad \varphi^2(x_{0,r}) = (\log r)^2. \]
		A rational model function restricts to a piecewise affine function of $\log r$ with finitely many pieces,
		whereas $(\log r)^2$ does not.
		For the piecewise affine behavior on skeleta, see \href{https://gubler.app.uni-regensburg.de/Publikationen/1404.7044v1.pdf}{Gubler--Rabinoff--Werner}.
	\end{latexmathchange}
\end{example}

\begin{proposition}[label=\labelNFM{prop:basic-properties-of-model-function}, status=draft, tag=0000000]
	Let $X$ be a variety over $K$.
	\begin{enumerate}
		\item $af+bg$ are model for \Yang{}
		\item $\max{f,g}$ and $\min{f,g}$ are \Yang{}
		\item model function is dense in $\scrC(X)$.
	\end{enumerate}
\end{proposition}
\begin{proof}[References for proof]
	\Yang{}
\end{proof}

\begin{remark}[label=\labelNFM{rmk:tropical-semiring-and-model-function}, status=draft, tag=0000000]
	\latexdeleted{The model functions usual does not form a ring since it cannot take multiplication.}
	\latexdeleted{but To Ai: fill this remark}
	\begin{latexmathchange}
		Model functions form a $\bbQ$-vector space and are closed under pointwise maximum and minimum.
		For sums and rational multiples, pass to a common model and add or scale the vertical divisors.
		To form a maximum, after passing to a common model, blow up the sum of the corresponding vertical fractional ideals so that it becomes invertible.
		The resulting model function is the pointwise maximum; the minimum follows by taking negatives.

		Define tropical operations by
		\[ \varphi \oplus \psi \coloneqq \max\{\varphi,\psi\}, \qquad \varphi \odot \psi \coloneqq \varphi + \psi. \]
		After adjoining the constant function $-\infty$, these operations make $\calM_v(X) \cup \{-\infty\}$ an idempotent semiring,
		with tropical zero $-\infty$ and tropical unit the ordinary zero function.
		Ordinary pointwise products need not be model functions, as $\varphi^2$ in the preceding example shows.
	\end{latexmathchange}
\end{remark}


\begin{definition}[label=\labelNFM{def:Chambert-Loir-measure}, status=draft, tag=0000000]
	Let $X$ and
	Let $\bar L$ be nef, with $L$ big.
	At every place $v$, there is a positive measure $c_1(\bar L_v)^d$ on $X_v^\an$ of total mass $\deg_L(X)$.
	At an infinite place it is the Bedford--Taylor product of the curvature current.
	At a finite place it is the Chambert-Loir measure: for model metrics it is obtained from intersection degrees on the special fiber, and semipositive metrics are treated by weak limits.
	We set
	\[ \mu_{\bar L,v} \coloneqq \frac{c_1(\bar L_v)^d}{\deg_L(X)}. \]
	Thus $\mu_{\bar L,v}$ is a probability measure.
	\Yang{}
\end{definition}




























\subsection[status=draft,tag=0000000,label=\labelNFM{subsec:adelic-line-bundle-and-positivity}]{Adelic line bundle and arithmetic intersection}

On this subsection, fix a number field $K$.
\Yang{}



\begin{definition}[label=\labelNFM{def:adelic-line-bundle-projective-case}, status=draft, tag=0000000]
	Let $X$ be a projective variety over $K$.
	An \emph{adelic line bundle} on $X$ is a pair
	\[ \bar L = (L,(\|\cdot\|_v)_{v \in M_K}) \]
	with continuous $v$-adic metrics, such that there is an nonempty open subset $U \subset \Spec \calO_K$ and an integral model $(\calX,\calL)/U$ of $(X,L)$ such that for all $v \in U$, $\|\cdot\|_v = \|\cdot\|_{\calX|_{\calO_v},\calL|_{\calO_v}}$ is the model metric induced by the pullback of $(\calX,\calL)$.
	\latexdeleted{It forms a group..}
	\latexmath{Their isometry classes form $\aPic(X)$ under tensor product, with duals as inverses.}
	\Yang{}
\end{definition}

\begin{definition}[label=\labelNFM{def:positivity-of-adelic-line-bundles}, status=draft, tag=0000000]
	An adelic line bundle is \emph{semipositive}, or \emph{relatively nef}, if all its local metrics are semipositive.
	It is \emph{nef} if it is semipositive and $h_{\bar L}(x) \geq 0$ for every $x \in X(\bar K)$.
	It is \emph{integrable} if $\bar L \cong \bar A \otimes \bar B^\vee$ for two nef adelic $\bbQ$-line bundles $\bar A,\bar B$.
	We denote the group of integrable adelic line bundles by $\aPic(X)_{\mathrm{int}}$.
	\Yang{}
\end{definition}

\begin{proposition}[label=\labelNFM{prop:fundamental-properties-of-adelic-line-bundle}, status=draft, tag=0000000]
	\latexdeleted{To Ai: finish}
	\begin{latexmathchange}
	Let $X$ be a projective variety over $K$.
	The following operations on line bundles extend to adelic line bundles.
	\begin{enumerate}
		\item Tensor products and duals carry the product and dual norms.
		      A morphism $g:Y \to X$ gives a pullback $g^*:\aPic(X) \to \aPic(Y)$.
		      A finite extension $E/K$ gives base change, with the metrics extended at every $w \mid v$.
		\item Tensor products, positive rational powers, and pullbacks preserve semipositivity and nefness.
		      Differences of nef $\bbQ$-line bundles are stable under tensor product, duals, and pullback, so the integrable classes form a $\bbQ$-vector space.
		\item For the point height defined below,
		      \[ h_{\bar L\otimes\bar M} = h_{\bar L}+h_{\bar M}, \qquad h_{g^*\bar L}(y) = h_{\bar L}(g(y)), \qquad h_{\bar L_E}(x) = h_{\bar L}(x). \]
		      If two metrics on the same $L$ satisfy $\|\cdot\|'_v = e^{-\varphi_v}\|\cdot\|_v$, then
		      \[ h_{\bar L'}(x)-h_{\bar L}(x) = \sum_{v \in M_K} n_v\int\varphi_v\,\upd\delta_{x,v}, \qquad |h_{\bar L'}(x)-h_{\bar L}(x)| \leq \sum_v n_v\|\varphi_v\|_\infty. \]
		      Here $\delta_{x,v}$ is the local Galois-orbit probability measure defined in the equidistribution subsection.
		      Only finitely many $\varphi_v$ are nonzero.
		\item If $L$ is ample, $h_{\bar L}$ is bounded below and has Northcott's property:
		      for every real $B$ and integer $D$, the set
		      \[ \{x \in X(\bar K):h_{\bar L}(x) \leq B,\ [K(x):K] \leq D\} \]
		      is finite.
	\end{enumerate}
	See \cite[Appendices~A.1.1--A.1.2 and A.5.2]{Yuan-Zhang-2026-AdelicLineBundles}.
	\end{latexmathchange}
\end{proposition}

\begin{proof}
	\leavevmode\par
	\begin{latexmathchange}
	Take a common model over an open subset of $\Spec\calO_K$ to check the adelic condition for the operations in (1).
	Local approximation by relatively nef models proves the semipositivity assertions; the height identities prove preservation of nonnegative point heights.
	The identities in (3) follow by adding logarithms of norms and by the local-degree formula under finite extensions.
	In particular, two adelic metrics on $L$ give heights differing by a bounded function.
	For $L$ ample, choose $m > 0$ and an embedding $\iota:X \hookrightarrow \bbP^N_K$ with $L^{\otimes m} \cong \iota^*\calO(1)$.
	Then $m h_{\bar L} = h_{\overline{\calO(1)}_{\max}}\circ\iota+O(1)$.
	The usual projective Northcott theorem and nonnegativity of the standard height give (4).
	\end{latexmathchange}
\end{proof}

\begin{definition}[label=\labelNFM{def:height-of-an-algebraic-point}, status=draft, tag=0000000]
	Let $x \in X(\bar K)$ and choose a finite extension $E/K$ with $x \in X(E)$.
	For a rational section $s$ of $L$ regular and nonzero at $x$, define
	\[ h_{\bar L}(x) \coloneqq -\sum_{w \in M_E}n_w\log\|s(x)\|_w, \qquad n_w = \frac{[E_w:\bbQ_p]}{[E:\bbQ]} \quad (w \mid p). \]
\end{definition}

\Yang{}
The metrics are pulled back from $K$.
\Yang{}
Only finitely many summands are nonzero.
Replacing $s(x)$ by $a s(x)$ changes the sum by $-\sum_w n_w\log|a|_w = 0$.
The local degree formula shows that the value is also independent of $E$.
Thus this gives a well-defined real-valued height on $X(\bar K)$.
\Yang{}

\begin{example}[label=\labelNFM{eg:standard-logarithmic-height}, status=draft, tag=0000000]
	Give $\scrO(1)$ on $\bbP^d$ the standard metric at every place, and denote the result by $\overline{\scrO(1)}_{\max}$.
	Applying the definition with a nonzero coordinate section gives
	\[ h_{\overline{\scrO(1)}_{\max}}([a_0:\cdots:a_d]) = \sum_{w \in M_E}n_w\log\max_i|a_i|_w. \]
	The product formula removes the term involving that coordinate.
	In particular, for $x = [x_0:\cdots:x_d] \in \bbP^d(\bbQ)$ with $\gcd(\{x_i\}) = 1$, we have
	\[  h_{\overline{\scrO(1)}_{\max}}(x) = \log \max\{|x_0|,\cdots,|x_d|\}. \]
\end{example}

\begin{definition}[label=\labelNFM{def:arithmetic-intersection-and-subvariety-height}, status=draft, tag=0000000]
	\latexdeleted{Let $\bar L_0,\ldots,\bar L_k$ be adelic line bundle on $X$ and $Z \subset X$ a closed subset variety of dimension $k$.}
	\latexdeleted{We define the \emph{arithmetic intersection number} as}
	\[ \latexdeletedmath{\widehat c_1(\bar L_0)\cdots\widehat c_1(\bar L_k)\cdot[Z] \coloneqq} \]
	\begin{latexmathchange}
	Let $\bar L_0,\ldots,\bar L_r$ be integrable adelic $\bbQ$-line bundles and $Z \subset X$ an integral subvariety of dimension $r$.
	The adelic Deligne pairing gives a metrized line over $\Spec K$,
	\[ \widehat c_1(\bar L_0)\cdots\widehat c_1(\bar L_r)\cdot[Z] \coloneqq \langle\bar L_0|_Z,\ldots,\bar L_r|_Z\rangle_{Z/K}. \]
	Its arithmetic degree is the real \emph{arithmetic intersection number}; we write
	\[ I_K(\bar L_0,\ldots,\bar L_r;Z) \coloneqq \adeg\bigl(\widehat c_1(\bar L_0)\cdots\widehat c_1(\bar L_r)\cdot[Z]\bigr). \]
	For model metrics, compute this degree by the Gillet--Soul\'e intersection on a common proper arithmetic model, dividing by the powers used to extend the $L_i$.
	For nef adelic bundles, take the limit of such model intersections; multilinearity then defines the value for differences of nef bundles.
	The limit is independent of the approximations.
	We use the unnormalized arithmetic degree: on $\Spec K$, it is $-\sum_v[K_v:\bbQ_p]\log\|s\|_v$ for $0 \neq s$ and $v \mid p$.
	See \cite[\S 4.1 and Appendix~A.5.3]{Yuan-Zhang-2026-AdelicLineBundles}.
	\end{latexmathchange}
	Let $Y \subset X$ be an integral closed subvariety of dimension $r$, and assume that $\deg_L(Y) = (L^r \cdot [Y]) > 0$.
	Its normalized height is
	\[ h_{\bar L}(Y) \coloneqq \frac{\adeg\bigl(\widehat c_1(\bar L)^{r+1}\cdot[Y]\bigr)}{(r+1)[K:\bbQ]\deg_L(Y)}. \]
	\Yang{}
\end{definition}

For a closed point $P$ with residue field $E$, one has $\deg_L(P) = [E:K]$, and this agrees with the preceding point height.
In particular,
\[ h_{\bar L}(X) = \frac{\adeg\bigl(\widehat c_1(\bar L)^{d+1}\cdot[X]\bigr)}{(d+1)[K:\bbQ]\deg_L(X)}. \]
See \cite[\S 3.2]{Chambert-Loir-2012-ArakelovGeometryHeights} and \cite[Appendix~A.5.3]{Yuan-Zhang-2026-AdelicLineBundles}.

\begin{theorem}[label=\labelNFM{thm:fundamental-properties-of-arithmetic-intersection}, status=draft, tag=0000000]
	\latexdeleted{Let ...}
	\latexdeleted{To Ai: finish}
	\begin{latexmathchange}
	For integrable adelic $\bbQ$-line bundles, the arithmetic intersection has the following properties.
	\begin{enumerate}
		\item It is symmetric and multilinear in the line bundles and additive in the cycle.
		      It is continuous for uniform convergence of semipositive metrics with fixed underlying bundles and fixed metrics outside a finite set of places.
		\item If $g:Y \to X$ is a morphism of projective varieties and $W \subset Y$ is an integral $r$-dimensional subvariety, the projection formula is
		      \[ I_K(g^*\bar L_0,\ldots,g^*\bar L_r;W) = I_K(\bar L_0,\ldots,\bar L_r;g_*[W]). \]
		      Here $g_*[W] = [K(W):K(g(W))][g(W)]$ when $\dim g(W) = r$, and is zero otherwise.
		\item For a finite extension $E/K$, with base-changed cycles interpreted additively,
		      \[ I_E(\bar L_{0,E},\ldots,\bar L_{r,E};Z_E) = [E:K]I_K(\bar L_0,\ldots,\bar L_r;Z). \]
		      Thus normalized heights are unchanged by base extension.
		\item Let $r = \dim X = d$, and change only the metric of $\bar L_0$ at $v$ by $\|\cdot\|_v' = e^{-\varphi}\|\cdot\|_v$, with the changed bundle integrable.
		      Then
		      \[ I_K(\bar L_0',\bar L_1,\ldots,\bar L_d;X)-I_K(\bar L_0,\bar L_1,\ldots,\bar L_d;X) = [K:\bbQ]n_v\int_{X_{\bar v}^\an}\varphi\,c_1(\bar L_{1,v})\cdots c_1(\bar L_{d,v}). \]
		\item If all the bundles are nef, their intersection on an effective cycle is nonnegative.
	\end{enumerate}
	See \cite[\S\S 4.1--4.4 and Appendix~A.5.3]{Yuan-Zhang-2026-AdelicLineBundles}.
	\end{latexmathchange}
\end{theorem}

\begin{proof}
	\leavevmode\par
	\begin{latexmathchange}
	For model metrics these are the symmetry, multilinearity, projection formula, metric-change formula, and positivity of arithmetic intersection.
	When the other bundles are semipositive, the metric-change formula bounds a variation by the uniform norm of the metric change times the relevant geometric degree.
	It therefore gives continuity and permits passage to the defining limits of nef adelic bundles.
	Expanding differences of nef bundles gives the integrable case; base change follows either on models or from the local-degree formula for the metrized Deligne pairing.
	\end{latexmathchange}
\end{proof}

\begin{definition}[label=\labelNFM{def:arithmetic-volume-of-adelic-line-bundle}, status=draft, tag=0000000]
	\latexdeleted{Let $X$ and $\bar{L}$ a ...}
	\[ \latexdeletedmath{a\vol(X) \coloneqq} \]
	\latexdeleted{To Ai: finish}
	\begin{latexmathchange}
	Let $X/K$ be projective of dimension $d$ and $\bar L$ an adelic line bundle.
	The set of \emph{small sections} and its logarithmic cardinality are
	\[ \widehat H^0(X,m\bar L) \coloneqq \{s \in H^0(X,L^{\otimes m}):\|s\|_{v,\sup} \leq 1\text{ for every }v\}, \qquad \widehat h^0(X,m\bar L) \coloneqq \log\#\widehat H^0(X,m\bar L). \]
	This is a finite set containing the zero section, rather than a $K$-vector space.
	The \emph{arithmetic volume} is
	\[ \widehat{\vol}(\bar L) \coloneqq \limsup_{m \to \infty}\frac{\widehat h^0(X,m\bar L)}{m^{d+1}/(d+1)!}. \]
	For $\bbQ$-line bundles, clear denominators and use homogeneity of degree $d+1$.
	The bundle is \emph{arithmetically big} if $\widehat{\vol}(\bar L) > 0$.
	This is stronger than bigness of the underlying $L$, which concerns growth of $\dim_K H^0(X,mL)$ in degree $d$.
	Our volume has no factor $1/[K:\bbQ]$.
	For a nef $\bar L$, the arithmetic Hilbert--Samuel formula gives
	\[ \widehat{\vol}(\bar L) = I_K(\bar L,\ldots,\bar L;X). \]
	See \cite[\S\S 5.1--5.2, especially Theorem~5.2.2]{Yuan-Zhang-2026-AdelicLineBundles}.
	\end{latexmathchange}
\end{definition}

\begin{theorem}[label=\labelNFM{thm:arithmetic-siu-inequality}, status=draft, tag=0000000]
	\latexdeleted{Let $X$ be}
	\Yang{}
	\latexdeleted{To Ai: finish}
	\begin{latexmathchange}
	Let $X/K$ be projective of dimension $d$, and let $\bar A,\bar B$ be nef adelic $\bbQ$-line bundles.
	Then the \emph{arithmetic Siu inequality} is
	\[ \widehat{\vol}(\bar A\otimes\bar B^\vee) \geq I_K(\bar A^{d+1};X)-(d+1)I_K(\bar A^d,\bar B;X). \]
	Here powers inside $I_K$ mean repeated arguments.
	In particular, a positive right-hand side implies that $\bar A\otimes\bar B^\vee$ is arithmetically big.
	This estimates the small sections left after subtracting a nef bundle, just as the geometric Siu inequality estimates ordinary sections.
	For the theorem and proof, see \cite[Theorem~5.2.2(2)]{Yuan-Zhang-2026-AdelicLineBundles}.
	\end{latexmathchange}
\end{theorem}


\begin{definition}[label=\labelNFM{def:successive-minima}, status=draft, tag=0000000]
	\latexdeleted{Let $X$ and $\bar{L}$ a ...}
	\[ \latexdeletedmath{e_i(\bar L) \coloneqq} \]
	\Yang{}
	\latexdeleted{To Ai: finish}
	\begin{latexmathchange}
	Let $\dim X = d$ and $\bar L$ be an adelic line bundle with $L$ ample.
	For $1 \leq i \leq d+1$, define the \emph{successive minima} by
	\[ e_i(\bar L) \coloneqq \sup_{\substack{Z \subset X_{\bar K}\text{ closed}\\\dim Z \leq d-i}}\ \inf_{x \in X(\bar K)\setminus Z(\bar K)} h_{\bar L}(x). \]
	We allow $Z = \varnothing$, with dimension $-\infty$.
	Then $e_1 \geq \cdots \geq e_{d+1}$; $e_1$ is the \emph{essential minimum}, and $e_{d+1} = \inf_x h_{\bar L}(x)$ is the absolute minimum.
	Equivalently,
	\[ e_1(\bar L) = \inf\{t \in \bbR:\{x \in X(\bar K):h_{\bar L}(x) \leq t\}\text{ is Zariski dense}\}. \]
	For every generic sequence $(x_m)$, one has $\liminf_m h_{\bar L}(x_m) \geq e_1(\bar L)$.
	Indeed, it eventually avoids each fixed proper closed $Z$ in the definition.
	The same definition of $e_1$ also makes sense when $L$ is only big.
	See \cite[Appendix~A.5.5]{Yuan-Zhang-2026-AdelicLineBundles}; that appendix uses arithmetic dimension $d+1$ in place of our geometric dimension $d$.
	\end{latexmathchange}
\end{definition}

\begin{theorem}[label=\labelNFM{thm:fundamental-inequality-IV}, status=draft, tag=0000000]
	\latexdeleted{Let $X$ be}

	\latexdeleted{To Ai: finish}
	\Yang{}
	\begin{latexmathchange}
	Let $X/K$ be projective of dimension $d$, let $L$ be ample, and let $\bar L$ be semipositive.
	Zhang's theorem of successive minima states
	\[ e_1(\bar L) \geq h_{\bar L}(X) \geq \frac{e_1(\bar L)+\cdots+e_{d+1}(\bar L)}{d+1}. \]
	If $\bar L$ is nef with big underlying $L$, one still has the fundamental inequalities
	\[ e_1(\bar L) \geq h_{\bar L}(X) \geq \frac{e_1(\bar L)}{d+1}. \]
	For the first statement and proof, see \cite[Theorem~A.5.2]{Yuan-Zhang-2026-AdelicLineBundles}; for the nef and big version, see Theorem~5.3.3 there.
	\end{latexmathchange}
\end{theorem}

\begin{latexmathchange}
	These inequalities explain the smallness condition in equidistribution.
	Every generic sequence has $\liminf h_{\bar L}(x_m) \geq e_1(\bar L) \geq h_{\bar L}(X)$.
	Thus a generic sequence whose heights tend to $h_{\bar L}(X)$ attains this lower bound.
	For a nef bundle, all the point heights are nonnegative, so the lower successive minima are nonnegative as well.
\end{latexmathchange}

































\subsection[status=draft,tag=0000000,label=\labelNFM{subsec:canonical-height}]{Canonical height}

\Yang{}
\latexdeleted{To Ai: finish this section}

On this subsection, let us turn to the dynamics.

\latexdeleted{Let $f: X \to X$ over ..}
\begin{latexmathchange}
	Let $X/K$ be a geometrically integral projective variety of dimension $d$ over a number field.
	A \emph{$q$-polarized dynamics} consists of $f:X \to X$, an ample line bundle $L$, and an isomorphism $\alpha:f^*L\xrightarrow{ \sim }L^{\otimes q}$, where $q \geq 2$.
	The morphism $f$ is finite and surjective, of degree $q^d$.
	Indeed, ampleness rules out a positive-dimensional fiber, and the degree follows from $(f^*L)^d = q^d L^d$.
\end{latexmathchange}

\begin{definition}[label=\labelNFM{def:canonical-height}, status=draft, tag=0000000]
	Let $(f,X)$ be a $q$-polarized dynamics with polarization $L$.
	\latexdeleted{Then there exists a canonical adelic line bundle structure $\bar{L}$ on $L$ such that $f^*\bar{L} = \bar{L}^{\otimes q}$.}
	\begin{latexmathchange}
	For the fixed polarization isomorphism $\alpha$, the \emph{canonical adelic line bundle} $\bar L_f$ is the unique adelic metric on $L$ for which $\alpha$ is an isometry
	\[ f^*\bar L_f \xrightarrow[\alpha]{ \sim } \bar L_f^{\otimes q}. \]
	Its point height $\widehat h_{f,L} \coloneqq h_{\bar L_f}$ is the \emph{canonical height}.
	The choice of $\alpha$ matters for the individual local metrics.
	Replacing $\alpha$ by $a\alpha$, with $a \in K^*$, changes these metrics by constants whose weighted sum is zero by the product formula; the point height is unchanged.
	\end{latexmathchange}
\end{definition}

\Needspace{22\baselineskip}
\begin{construction}[label=\labelNFM{constr:canonical-metric-by-iteration}, status=draft, tag=0000000]
	\leavevmode\par
	\begin{latexmathchange}
	Start with a semipositive adelic metric $\bar L_0$ and use the iterates of $\alpha$ to regard
	\[ \bar L_n \coloneqq q^{-n}(f^n)^*\bar L_0 \]
	as metrics on the same $L$; multiplication by $q^{-n}$ means taking the $q^n$-th root of the norms.
	At each place, let $C_v = \|\log(\|\cdot\|_{1,v}/\|\cdot\|_{0,v})\|_\infty$.
	The pullback operation divided by $q$ contracts differences of logarithmic metrics, so
	\[ \left\|\log\frac{\|\cdot\|_{n+1,v}}{\|\cdot\|_{n,v}}\right\|_\infty \leq q^{-n}C_v. \]
	The geometric series converges, giving a continuous limiting metric at every place.
	Choose a model on which $f,L,\alpha$ all extend away from finitely many places; there $C_v = 0$.
	Hence the limit is adelic, and it is semipositive because the iterated metrics are semipositive.
	It satisfies $f^*\bar L_f \cong \bar L_f^{\otimes q}$ isometrically.
	For two metrics with this property, their logarithmic difference $u_v$ satisfies $u_v = q^{-1}u_v\circ f$.
	The supremum norm forces $u_v = 0$, proving uniqueness and independence of the starting metric.
	See \cite[Theorem~6.1.1 and Appendix~A.5.2]{Yuan-Zhang-2026-AdelicLineBundles}.
	\end{latexmathchange}
\end{construction}

\begin{proposition}[label=\labelNFM{prop:canonical-height-properties}, status=draft, tag=0000000]
	\leavevmode\par
	\begin{latexmathchange}
	For any Weil height $h_L$ associated to $L$, one has
	\[ \widehat h_{f,L}(x) = \lim_{n \to \infty}q^{-n}h_L(f^n(x)), \qquad \widehat h_{f,L}(f(x)) = q\widehat h_{f,L}(x), \qquad \widehat h_{f,L} = h_L+O(1). \]
	Moreover, $\widehat h_{f,L} \geq 0$, and
	\[ \widehat h_{f,L}(x) = 0 \quad \Longleftrightarrow \quad x\text{ is preperiodic for }f. \]
	The canonical bundle is nef and $h_{\bar L_f}(X) = 0$.
	In particular, periodic points have height zero, and for $y \in f^{-n}(a)$,
	\[ \widehat h_{f,L}(y) = q^{-n}\widehat h_{f,L}(a). \]
	\end{latexmathchange}
\end{proposition}
\begin{proof}
	\leavevmode\par
	\begin{latexmathchange}
	The height machine gives $|h_L(f(x))-qh_L(x)| \leq C$ for a constant $C$ independent of $x$.
	The successive differences in $q^{-n}h_L(f^n(x))$ are at most $Cq^{-n-1}$.
	Thus the sequence converges, differs from $h_L$ by at most $C/(q-1)$, and satisfies the functional equation.
	The metric construction gives the same limit by the metric-change bound above.
	An ample height is bounded below, so dividing that lower bound by $q^n$ shows $\widehat h_{f,L}(x) \geq 0$.

	A finite orbit has canonical height zero by the functional equation.
	Conversely, if $\widehat h_{f,L}(x) = 0$, all iterates have bounded $h_L$ and belong to the fixed number field $K(x)$.
	Northcott's theorem makes the orbit finite.
	Semipositivity and nonnegative point heights show that $\bar L_f$ is nef.
	Finally the projection formula gives
	\[ q^{d+1}I_K(\bar L_f^{d+1};X) = I_K((f^*\bar L_f)^{d+1};X) = q^d I_K(\bar L_f^{d+1};X). \]
	Since $q > 1$, this intersection is zero, so its normalized height is zero.
	The backward-orbit formula follows from the functional equation.
	\end{latexmathchange}
\end{proof}

\begin{remark}[label=\labelNFM{rmk:canonical-metric-and-green-current}, status=draft, tag=0000000]
	\leavevmode\par
	\begin{latexmathchange}
	For $X = \bbP^d$ and $L = \calO(1)$, start at an infinite place with the Fubini--Study metric.
	The curvatures of its iterates converge to
	\[ c_1(\bar L_{f,v}) = \lim_{n \to \infty}q^{-n}(f^n)^*\omega_\FS = T_f, \qquad \mu_{\bar L_f,v} = T_f^d = \mu_f. \]
	Thus the analytic Green current and the arithmetic canonical height come from the same metrized line bundle.
	The canonical height of $X$ is zero, so the smallness condition in arithmetic equidistribution is exactly $\widehat h_{f,L}(x_m) \to 0$.
	Genericity is a separate condition: small points could all lie in a fixed proper invariant subvariety.
	See \cite[Appendix~A.6.4]{Yuan-Zhang-2026-AdelicLineBundles}.
	\end{latexmathchange}
\end{remark}





